\documentclass[10pt]{article}
\usepackage[T1]{fontenc}
\usepackage[utf8]{inputenc}
\usepackage[french]{babel}
\usepackage{amsmath,amssymb,mathrsfs,mathpazo}
\usepackage{geometry}
\geometry{a4paper,margin=2cm}
\usepackage{graphicx,array,tabularx,tabularray,colortbl,multirow,multicol,booktabs,enumitem,pifont,ulem,fancybox,diagbox,cancel}
\usepackage{pstricks,pstricks-add,pst-plot,pst-tree,pst-node,pst-text,pst-func,pst-math,pst-bspline,pst-3dplot,pst-eucl,pst-solides3d}
\let\pscalculate\undefined
\usepackage{pst-bezier}
\usepackage{tikz,pgfplots,tkz-tab}
\usetikzlibrary{arrows,arrows.meta,patterns,calc,positioning,trees,decorations.markings,shapes.geometric,shapes.misc,shapes.arrows,intersections,decorations.pathmorphing,decorations.pathreplacing,backgrounds,shadows,babel}
\pgfplotsset{compat=1.18}
\usepgfplotslibrary{fillbetween}
\usepackage[np]{numprint}
\usepackage{listings,makecell,pas-tableur,twoopt}
\let\barre\undefined
\usepackage[pstricks]{tablvar}
\definecolor{darkgreen}{RGB}{0,100,0}
\providecommand{\red}{\color{red}}
\providecommand{\blue}{\color{blue}}
\providecommand{\green}{\color{green}}
\providecommand{\black}{\color{black}}
\providecommand{\white}{\color{white}}
\providecommand{\gray}{\color{gray}}
\providecommand{\cyan}{\color{cyan}}
\providecommand{\magenta}{\color{magenta}}
\providecommand{\yellow}{\color{yellow}}
\providecommand{\violet}{\color{violet}}
\providecommand{\orange}{\color{orange}}
\providecommand{\euro}{\texteuro}
\providecommand{\EUR}{\texteuro}
\providecommand{\ieme}{\textsuperscript{e}}
\providecommand{\degres}{\ensuremath{^\circ}}
\providecommand{\degree}{\ensuremath{^\circ}}
\providecommand{\vv}[1]{\overrightarrow{#1}}
\providecommand{\Coord}[1]{#1}
\providecommand{\decosix}{}
\providecommand{\starredbullet}{\ensuremath{\bullet}}
\pagestyle{empty}
\setlength{\parindent}{0pt}
\usepackage{adjustbox}\DeclareUnicodeCharacter{2236}{:}\DeclareUnicodeCharacter{2190}{\ensuremath{\leftarrow}}
\providecommand{\R}{\mathbb{R}}
\providecommand{\N}{\mathbb{N}}
\providecommand{\Z}{\mathbb{Z}}
\providecommand{\Q}{\mathbb{Q}}
\providecommand{\D}{\mathbb{D}}
\providecommand{\C}{\mathbb{C}}
\providecommand{\e}{{\mathrm{e}}}
\providecommand{\vect}[1]{\overrightarrow{#1}}
\providecommand{\vectt}[1]{\overrightarrow{\mathrm{#1}}}
\providecommand{\vvt}[1]{\overrightarrow{\mathrm{#1}}}
\providecommand{\np}[2][]{{#2~#1}}
\providecommand{\vv}[1]{{\overrightarrow{#1}}}
\renewcommand{\d}{{\mathrm{d}}}
\providecommand{\degres}{{^\circ}}
\providecommand{\degre}{{^\circ}}
\providecommand{\blue}{\color{blue}}
\providecommand{\red}{\color{red}}
\providecommand{\enskip}{\quad}
\providecommand{\llbracket}{\lbrack\!\lbrack}
\providecommand{\rrbracket}{\rbrack\!\rbrack}
\providecommand{\euro}{{\text{€}}}
\providecommand{\no}{{\text{n°}}}
\providecommand{\interval}[2]{{[#1;#2]}}
\providecommand{\intervFO}[2]{{[#1;#2[}}
\providecommand{\dotfill}{\ldots}

\providecommand{\ds}{\displaystyle}
\providecommand{\textding}[1]{#1}
\providecommand{\Oij}{\(\mathrm{O};\vec{\imath},\vec{\jmath}\)}
\providecommand{\Oijk}{\(\mathrm{O};\vec{\imath},\vec{\jmath},\vec{k}\)}
\providecommand{\Ouv}{\(\mathrm{O};\vec{u},\vec{v}\)}
\providecommand{\intervFF}[2]{\ensuremath{[#1;#2]}}
\renewcommand{\euro}{\ifmmode\text{\texteuro}\else\texteuro\fi}
\DeclareUnicodeCharacter{00B0}{\ifmmode^{\circ}\else\textdegree\fi}
\DeclareUnicodeCharacter{20AC}{\ifmmode\text{\texteuro}\else\texteuro\fi}
% Énoncés : Éducation nationale. Transcription et corrigés : APMEP, Denis Vergès et contributeurs. Découpage et classement : CoopMaths, Sébastien Lozano, Éric Elter et contributeurs. Adaptation de la présentation : Maths Explorer.
% CC BY-SA : https://coopmaths.fr/www/about/
% https://coopmaths.fr/alea/static/bac/2025/tex/bac_2025_05_sujet1_ameriquenord_2.tex
\begin{document}
\noindent\textit{APMEP / CoopMaths. CC BY-SA. Présentation adaptée par Maths Explorer.}\par\medskip
\medskip

On considère la suite numérique $\left(u_{n}\right)$ définie par son premier terme $u_{0}=2$ et pour tout entier naturel $n$, par :


\[u_{n+1}=\dfrac{2 u_{n}+1}{u_{n}+2}\]

On admet que la suite $\left(u_{n}\right)$ est bien définie.

\begin{enumerate}
\item Calculer le terme $u_{1}$.
\item On définit la suite $\left(a_{n}\right)$ pour tout entier naturel $n$, par :

\[a_{n}=\dfrac{u_{n}}{u_{n}-1}\]

On admet que la suite $\left(a_{n}\right)$ est bien définie.

	\begin{enumerate}
		\item Calculer $a_{0}$ et $a_{1}$.

		\item Démontrer que, pour tout entier naturel $n$, $a_{n+1}=3 a_{n}-1$.

		\item Démontrer par récurrence que, pour tout entier naturel $n$ supérieur ou égal à~1,

\[a_{n} \geqslant 3 n-1\]

		\item En déduire la limite de la suite $\left(a_{n}\right)$.
	\end{enumerate}
\item On souhaite étudier la limite de la suite $\left(u_{n}\right)$.
	\begin{enumerate}
		\item Démontrer que pour tout entier naturel $n$, $u_{n}=\dfrac{a_{n}}{a_{n}-1}$.

		\item En déduire la limite de la suite $\left(u_{n}\right)$.
		\end{enumerate}
\item On admet que la suite $\left(u_{n}\right)$ est décroissante.

On considère le programme suivant écrit en langage Python :

\begin{center}
\begin{ttfamily}
\begin{tabularx}{8cm}{|l @{\quad} X|}\hline
	1& \textbf{def} algo(p):\\
	2& \quad u=2\\
	3& \quad n=0\\
	4& \quad \textbf{while} u-1>p:\\
	5& \quad\quad u=(2*u+1)/(u+2)\\
	6& \quad\quad n=n+1\\
	7& \quad \textbf{return} (n,u)\\ \hline
\end{tabularx}
\end{ttfamily}
\end{center}

	\begin{enumerate}
		\item Interpréter les valeurs \texttt{n} et \texttt{u} renvoyées par l'appel de la fonction \texttt{algo(p)} dans le contexte de l'exercice.
		\item Donner, sans justifier, la valeur de $n$ pour $p=0,001$.
	\end{enumerate}
\end{enumerate}

\end{document}
